%%=============================================================================
%% PDF-pack Conceptblok A (levels 1-3)
%%
%% Standalone document dat aan deelnemers wordt overhandigd in de "traditionele"
%% conditie van de gebruikersstudie (zie hoofdstuk Evaluatie van de bachelorproef).
%%
%% Compileren: xelatex pdfpack-blokA.tex
%%=============================================================================

\documentclass[11pt,a4paper]{article}

\usepackage[dutch]{babel}
\usepackage[utf8]{inputenc}
\usepackage[T1]{fontenc}
\usepackage[margin=2.5cm]{geometry}
\usepackage{xcolor}
\usepackage{listings}
\usepackage{enumitem}
\usepackage{titlesec}
\usepackage{fancyhdr}
\usepackage{tabularx}
\usepackage{amssymb}
\usepackage{hyperref}

\definecolor{cobolblue}{RGB}{0,80,140}
\definecolor{lightgray}{RGB}{240,240,240}

\lstset{
    basicstyle=\ttfamily\small,
    backgroundcolor=\color{lightgray},
    frame=single,
    framesep=4pt,
    framerule=0pt,
    rulecolor=\color{lightgray},
    breaklines=true,
    columns=fullflexible,
    keywordstyle=\color{cobolblue}\bfseries,
    morekeywords={IDENTIFICATION,DIVISION,PROGRAM-ID,ENVIRONMENT,DATA,WORKING-STORAGE,
        SECTION,PROCEDURE,STOP,RUN,DISPLAY,PIC,PICTURE,VALUE,MOVE,IF,ELSE,END-IF,
        PERFORM,VARYING,UNTIL,END-PERFORM,COMPUTE,FILE,STATUS,SELECT,ASSIGN,FD,
        OPEN,READ,CLOSE,INPUT,OUTPUT}
}

\pagestyle{fancy}
\fancyhf{}
\fancyhead[L]{\textsc{Learn COBOL \textbar\ PDF-pack Conceptblok A}}
\fancyhead[R]{Pagina \thepage}
\fancyfoot[C]{Bachelorproef Silas De Vuyst, HOGENT, 2025--2026}

\titleformat{\section}{\Large\bfseries\color{cobolblue}}{\thesection}{1em}{}
\titleformat{\subsection}{\large\bfseries}{\thesubsection}{1em}{}

\hypersetup{colorlinks=true,linkcolor=cobolblue,urlcolor=cobolblue}

\begin{document}

%%---------- Titelpagina -----------------------------------------------------
\begin{titlepage}
\begin{center}
\vspace*{2cm}
{\Huge\bfseries Conceptblok A}\\[0.5cm]
{\Large Programmaskelet, PICTURE-clausules en datahiërarchie}\\[2cm]

{\large \textit{Traditioneel leermateriaal voor de gebruikersstudie}}\\[0.3cm]
{\large Levels 1 tot en met 3}\\[3cm]

\rule{\textwidth}{0.4pt}\\[0.5cm]
\begin{minipage}{0.9\textwidth}
\textbf{Voor de deelnemer:}
Lees dit pakket aandachtig door en los daarna de opdracht op pagina~\pageref{sec:opdracht-A} op. 
Je hebt 20~minuten om alle leerdoelen op de checklist (\,p.~\pageref{sec:rubric-A})\, te behalen. 
Schrijf je code op papier of in een editor naar keuze, zonder gebruik van automatische validatietools. 
Markeer telkens een nieuwe ``poging'' wanneer je je code aanpast en opnieuw wilt controleren.
\end{minipage}\\
\rule{\textwidth}{0.4pt}\\[1cm]

\vfill
\textbf{Bachelorproef Toegepaste Informatica} \\
Silas De Vuyst \textbullet\ HOGENT \textbullet\ academiejaar 2025--2026
\end{center}
\end{titlepage}

%%---------- Inhoud ----------------------------------------------------------
\section{Inleiding}

COBOL (\textit{Common Business Oriented Language}) is een procedurele programmeertaal uit 1960 die nog steeds aan de basis ligt van een groot deel van het wereldwijde betalings\-verkeer. 
In dit conceptblok maak je kennis met de drie meest fundamentele bouwstenen die elk COBOL-programma gemeen heeft:

\begin{enumerate}[label=(\arabic*)]
    \item de strakke programmastructuur in vier \textbf{divisions},
    \item de \textbf{\texttt{PICTURE}}-clausule die het opslagformaat van data vastlegt,
    \item de \textbf{datahiërarchie} via level numbers.
\end{enumerate}

Daar waar relevant wordt een korte vergelijking met Python gemaakt, zodat je de nieuwe constructies kan koppelen aan wat je al kent.

\section{Level 1: De vier divisions}

Een COBOL-programma is opgebouwd uit vier verplichte \textbf{divisions}, in een vaste volgorde:

\begin{description}[leftmargin=4cm,labelwidth=3.5cm,style=nextline]
    \item[\texttt{IDENTIFICATION DIVISION}] Metadata van het programma. Verplicht is enkel \texttt{PROGRAM-ID}, dat als unieke naam fungeert.
    \item[\texttt{ENVIRONMENT DIVISION}] Koppeling met de hardware/omgeving (bestanden, externe systemen). In eenvoudige programma's vaak leeg.
    \item[\texttt{DATA DIVISION}] Declaratie van alle variabelen, in secties zoals \texttt{WORKING-STORAGE SECTION} en \texttt{FILE SECTION}.
    \item[\texttt{PROCEDURE DIVISION}] De uitvoerbare logica, gestructureerd in \textit{paragraphs}. Eindigt met \texttt{STOP RUN.}
\end{description}

\textbf{Vergelijking met Python.} Een Python-script kent geen verplichte structuur: imports, functies en \texttt{main()} mogen door elkaar staan. In COBOL is de volgorde strikt en niet-onderhandelbaar.

\subsection*{Voorbeeld}
\begin{lstlisting}
       IDENTIFICATION DIVISION.
       PROGRAM-ID. HELLO.

       ENVIRONMENT DIVISION.

       DATA DIVISION.
       WORKING-STORAGE SECTION.

       PROCEDURE DIVISION.
           DISPLAY "OK".
           STOP RUN.
\end{lstlisting}

\section{Level 2: De PICTURE-clausule}

In Python is een variabele een verwijzing naar een object; het type wordt dynamisch bepaald. 
In COBOL is een variabele een reeks bytes, en de \textbf{\texttt{PICTURE}}-clausule (afgekort \texttt{PIC}) bepaalt hoe die bytes geïnterpreteerd worden.

De drie symbolen die je hier nodig hebt:

\begin{itemize}
    \item \textbf{\texttt{9}}: één numeriek cijfer (0--9). \texttt{PIC 9(03)} is dus een geheel getal van drie cijfers.
    \item \textbf{\texttt{X}}: één alfanumeriek karakter. \texttt{PIC X(10)} is een tekenreeks van tien karakters.
    \item \textbf{\texttt{V}}: \emph{impliciet} decimaalteken. \texttt{PIC 9(3)V99} reserveert drie cijfers voor het gehele deel en twee voor de decimalen, zonder dat de komma zelf een byte inneemt.
\end{itemize}

\subsection*{Voorbeeld}
\begin{lstlisting}
       01 KLANT.
          05 NAAM     PIC X(20)   VALUE "Jan Janssens".
          05 LEEFTIJD PIC 9(03)   VALUE 042.
          05 SALDO    PIC 9(5)V99 VALUE 12345.67.
\end{lstlisting}

\textbf{Initialisatie zonder \texttt{VALUE}.} Numerieke velden worden automatisch op nullen gezet, alfanumerieke velden op spaties.

\section{Level 3: Data-hiërarchie}

COBOL gebruikt \textbf{level numbers} om een hiërarchische datastructuur op te bouwen, vergelijkbaar met geneste \texttt{dataclass}'es of \texttt{dict}'s in Python.

\begin{itemize}
    \item \textbf{Level \texttt{01}}: definieert een record (de wortel).
    \item \textbf{Level \texttt{05}--\texttt{49}}: groepsitems en velden binnen het record.
    \item \textbf{Level \texttt{77}}: een losstaande variabele die niet bij een record hoort.
\end{itemize}

\subsection*{Voorbeeld: Python versus COBOL}

\begin{lstlisting}
# Python
employee = {
    "id": 12345,
    "name": {
        "first": "Jan",
        "last":  "Janssens"
    }
}
\end{lstlisting}

\begin{lstlisting}
       01 EMPLOYEE-DATA.
          05 EMP-ID    PIC 9(05).
          05 EMP-NAME.
             10 FIRST  PIC X(10).
             10 LAST   PIC X(10).
\end{lstlisting}

%%---------- Opdracht --------------------------------------------------------
\newpage
\section{Opdracht voor conceptblok A}
\label{sec:opdracht-A}

Schrijf één COBOL-programma met de naam \texttt{KLANT-INFO} dat \emph{alle} onderstaande leerdoelen tegelijkertijd realiseert. 
Je hoeft de code niet uit te voeren: de evaluatie gebeurt op basis van de checklist hieronder.

\begin{enumerate}
    \item Bevat alle vier verplichte divisions in de juiste volgorde.
    \item Eindigt de \texttt{PROCEDURE DIVISION} met \texttt{STOP RUN.}
    \item Declareert minstens één numeriek veld met \texttt{PIC 9(\ldots{})}.
    \item Declareert minstens één alfanumeriek veld met \texttt{PIC X(\ldots{})}.
    \item Declareert minstens één veld met een impliciete decimaal (\texttt{V}), bijvoorbeeld voor een bedrag.
    \item Bouwt minstens één hiërarchisch record met een \texttt{01}-niveau en minstens één geneste \texttt{05} (\textit{groep + subvelden}).
\end{enumerate}

\textbf{Praktisch.} Schrijf je oplossing op de volgende pagina (of in een editor naar keuze, mits je het resultaat printscreent of overneemt). 
Markeer iedere keer dat je je code controleert tegen de uitleg en aanpast als één ``poging'', en houd zelf bij hoeveel pogingen je nodig had.

%%---------- Rubric ----------------------------------------------------------
\newpage
\section{Zelfcontrole: rubric}
\label{sec:rubric-A}

Vink na het indienen aan welke leerdoelen je oplossing bereikt heeft. 
Een leerdoel telt enkel als ``geslaagd'' wanneer het zonder voorbehoud aanwezig is in jouw code.

\begin{tabularx}{\textwidth}{|c|X|}
\hline
\textbf{\#} & \textbf{Leerdoel} \\
\hline
$\square$ & 1. Alle vier divisions aanwezig en in de juiste volgorde. \\
\hline
$\square$ & 2. \texttt{PROCEDURE DIVISION} eindigt met \texttt{STOP RUN.} \\
\hline
$\square$ & 3. Minstens één veld met \texttt{PIC 9(\ldots{})}. \\
\hline
$\square$ & 4. Minstens één veld met \texttt{PIC X(\ldots{})}. \\
\hline
$\square$ & 5. Minstens één veld met impliciete decimaal (\texttt{V}). \\
\hline
$\square$ & 6. Minstens één \texttt{01}-record met geneste \texttt{05}-subvelden. \\
\hline
\end{tabularx}

\bigskip
\noindent\textbf{Aantal pogingen:}\ \rule{2cm}{0.4pt} \hfill \textbf{Tijd-tot-succes:}\ \rule{2cm}{0.4pt}~min

\bigskip
\noindent\textit{Dank voor je deelname.} Lever dit pakket en je oplossing in bij de onderzoeker.

\end{document}
